\documentclass[11pt,a4paper]{article}
\usepackage[margin=22mm,headheight=15pt]{geometry}
\usepackage[T1]{fontenc}
\usepackage{amsmath,amssymb,amsthm}
\usepackage{newtxtext,newtxmath}
\usepackage{microtype}
\usepackage{booktabs,array}
\usepackage[dvipsnames]{xcolor}
\usepackage{enumitem}
\usepackage{listings}
\usepackage{fancyhdr}
\usepackage{xurl}
\usepackage{hyperref}
\definecolor{inkblue}{HTML}{173650}
\definecolor{lightgray}{HTML}{F3F5F7}
\hypersetup{colorlinks=true,linkcolor=inkblue,citecolor=inkblue,urlcolor=inkblue,
  pdftitle={A quadratic-cost backward-error solver for general real systems},
  pdfauthor={Sidney Holden},
  pdfsubject={An affirmative solution to IE-12 in the stated exact-real model}}
\pagestyle{fancy}
\fancyhf{}
\fancyhead[L]{\small\scshape IE-12 \quad Exact-real backward-error complexity}
\fancyhead[R]{\small 12 September 2026}
\fancyfoot[C]{\small\thepage}
\renewcommand{\headrulewidth}{0.3pt}
\setlength{\parindent}{1.2em}
\setlength{\parskip}{3pt}
\setlength{\emergencystretch}{2em}
\newtheorem{theorem}{Theorem}
\newtheorem{lemma}[theorem]{Lemma}
\newtheorem{proposition}[theorem]{Proposition}
\theoremstyle{remark}
\newtheorem{remark}[theorem]{Remark}
\newcommand{\R}{\mathbb R}
\newcommand{\Z}{\mathbb Z}
\newcommand{\E}{\mathbb E}
\newcommand{\PP}{\mathbb P}
\newcommand{\eps}{\varepsilon}
\newcommand{\norm}[1]{\left\|#1\right\|_2}
\newcommand{\abs}[1]{\left|#1\right|}
\newcommand{\berr}{\operatorname{berr}}
\newcommand{\sgnp}{\operatorname{sign}_{+}}
\lstset{basicstyle=\small\ttfamily,columns=fullflexible,keepspaces=true,
  frame=single,rulecolor=\color{black!20},backgroundcolor=\color{lightgray},
  xleftmargin=5pt,xrightmargin=5pt,framesep=6pt,breaklines=true,
  showstringspaces=false,aboveskip=10pt,belowskip=10pt}

\begin{document}
\thispagestyle{plain}
\begin{center}
{\color{inkblue}\Large\bfseries A quadratic-cost backward-error solver\\[4pt]
for general real systems}\par
\vspace{7pt}
{\large An affirmative solution to IE-12 in the stated exact-real model}\par
\vspace{9pt}
{\large Sidney Holden}\par
{\small Center for Computational Biology, Flatiron Institute, Simons Foundation}\par
{12 September 2026}
\end{center}

\begin{abstract}
We give an explicit randomized algorithm for the normalized, fully accessible
real linear-system problem in IE-12. Its deterministic worst-case operation
count is $O(n^2\eps^{-3})$; with probability greater than $0.997$ its nonzero
output satisfies $\norm{Ax-b}/\norm{x}\le 5\eps/8$. Thus the answer to IE-12
is affirmative with $q=3$. The algorithm rounds $A$ once to a nearby scaled
integer matrix. A weighted-pattern multiplication scheme supports repeated
products with this matrix and its transpose more cheaply than dense products.
A rectangular augmented system supplies an exact kernel, and a fixed-length
randomly started polynomial filter finds an approximate kernel vector.
A clamping step converts it to the required backward-error guarantee without
perturbing $b$. The logarithm in the iteration count is canceled by the
preprocessed product cost, not removed from convergence itself. All essential
lemmas, probability estimates, model details, and cost bounds are proved below.
\end{abstract}

\noindent\textbf{Verification status.} This is a complete mathematical argument
with executable exact-arithmetic component tests and floating-point
illustrations. The submitted draft was prepared with ChatGPT assistance;
authorship is recorded as Sidney Holden at his request. Independent informal
Codex AI-agent review is documented in the accompanying submission record.
This is not external human peer review or proof-assistant verification. The theorem is about exact-real scalar complexity, not
finite-precision stability, bit complexity, or practical speed.

\section{Problem, result, and mechanism}

IE-12 asks for absolute constants $C,q>0$ and an algorithm that, given a
nonsingular $A\in\R^{n\times n}$ with promised $\norm{A}=1$, a nonzero
$b\in\R^n$, and $0<\eps<1/2$, uses at most $Cn^2\eps^{-q}$ operations and
returns a nonzero vector satisfying
\begin{equation}\label{eq:target}
 \PP\!\left\{\frac{\norm{Ax-b}}{\norm{A}\norm{x}}\le\eps\right\}\ge0.99.
\end{equation}
The allowed primitives are exact-real scalar arithmetic, square roots,
comparisons, and independent standard Gaussian or random-bit draws. Input
entries are accessible; normalization need not be computed. The metric allows
perturbations in $A$ only \cite{ie12}.

The cited general-system result has cost
$O(n^2\log(n/\delta)/\eps)$ for failure probability $\delta$; the corresponding
PSD result is $O(n^2/\sqrt\eps)$ \cite[Corollaries 17 and 25]{dnr}.
Our construction addresses the stated operation-count question by exploiting
entry access and preprocessing. It does not assert dimension-independent
convergence of a black-box Krylov iteration.

\begin{theorem}[Affirmative answer to IE-12]\label{thm:main}
There is an explicit randomized algorithm in the stated exact-real model and
an absolute constant $C$ such that, for every $n\ge1$, every real
$A\in\R^{n\times n}$ with $\norm{A}=1$, every $b\ne0$, and every
$0<\eps<1/2$, the algorithm always returns $x\ne0$, uses at most
$Cn^2\eps^{-3}$ operations and $O(n^2)$ real-scalar storage, and satisfies
\begin{equation}\label{eq:main}
 \PP\!\left\{\frac{\norm{Ax-b}}{\norm{x}}\le\frac{5\eps}{8}\right\}>0.997.
\end{equation}
Its output also always satisfies $\norm{x}\le8\norm{b}/\eps$.
The constant is independent of all input data and of conditioning.
Nonsingularity is not required for this stronger statement.
\end{theorem}

The matrix-product construction is related to the Mailman method, which saves
a logarithmic factor for preprocessed constant-alphabet matrices
\cite{mailman}. Here the rounded matrix has an alphabet that can grow with
$n$ and $1/\eps$, so applying that result as a constant-alphabet black box
would not suffice. We instead charge an integer coefficient $a$ a weight
$1+|a|$, enumerate only patterns of bounded total weight, and process heavy
entries individually. The resulting estimates are
\begin{align}
 W&=O(n^2/\eps),\label{eq:overviewW}\\
 \text{one product}&=O(n4^k+W/k+n),\label{eq:overviewproduct}\\
 \text{number of iterations}&=O(k/\eps^2),\label{eq:overviewiter}
\end{align}
where $k=\max\{1,\lfloor\log_{16}n\rfloor\}$. Since $k4^k=O(n)$,
multiplying \eqref{eq:overviewproduct} by \eqref{eq:overviewiter} gives
$O(n^2/\eps^3)$. Both preprocessing and every table computation are included.

Throughout, vector and matrix norms without a subscript other than $F$ are
Euclidean and spectral norms, respectively. Transpose is written $\top$.
Indexed scalar-array reads and writes can be counted at unit cost without
changing any bound. Equivalently, preprocessing builds a fixed-wiring scalar
arithmetic circuit for each subsequent product. No floor, logarithm, hashing,
SVD, or real-number digit-extraction oracle is assumed.

\section{One bounded-cost rounding of the input}

Set
\begin{equation}\label{eq:parameters}
 \beta=\norm{b},\qquad c=b/\beta,\qquad
 \rho=\eta=\eps/8,\qquad h=\frac{\eps}{64\sqrt n}
      =\frac{\rho}{8\sqrt n}.
\end{equation}
For each $t=A_{ij}/h$, let $m=\lfloor t\rfloor$ and $p=t-m$.
Independently set $Z_{ij}=m+1$ with probability $p$ and $Z_{ij}=m$ otherwise.
Write $Q=hZ$.

\subsection{Implementation using exactly the permitted primitives}

The floor is not taken as a primitive. Starting with $m=0$, if $t\ge0$,
increment $m$ until $m+1>t$; if $t<0$, decrement $m$ until $m\le t$.
This costs $O(1+|t|)$ scalar additions and comparisons and gives the exact floor.

An exact uniform random variable can be obtained in constant cost from three
independent standard Gaussians. Put
\begin{equation}\label{eq:uniform}
 R=\sqrt{G_1^2+G_2^2+G_3^2},\qquad
 U=\frac{1+G_1/R}{2}\quad(R>0).
\end{equation}
If $R=0$, assign $U=1/2$ without retrying. This event has probability zero.
The normalized Gaussian vector is uniform on the unit sphere in $\R^3$,
by rotational invariance. A strip whose first coordinate lies in
$[z,z+dz]$ has area $2\pi\,dz$, whereas the sphere has area $4\pi$.
Consequently $G_1/R$ is uniform on $[-1,1]$, and $U$ is uniform on $[0,1]$.
The comparison $U<p$ therefore implements the rounding rule exactly. Independent
triples are used for distinct entries, and the later starting vector uses fresh
draws. There is no rejection loop and no expected-time qualification.

\begin{lemma}[Weight, cost, and spectral error]\label{lem:round}
For every realization of the rounding,
\begin{equation}\label{eq:weight}
 W:=\sum_{i,j}(1+|Z_{ij}|)
 \le (2+64/\eps)n^2\le65n^2/\eps.
\end{equation}
The entire rounding costs $O(n^2/\eps)$ operations, including computation of the
integer parts. Moreover,
\begin{equation}\label{eq:roundprob}
 \PP\{\norm{Q-A}>\rho\}
 \le 2\exp\bigl((2\log 9-32)n\bigr)<10^{-3}.
\end{equation}
\end{lemma}

\begin{proof}
Since $|Z_{ij}|\le |A_{ij}|/h+1$, Cauchy--Schwarz gives
\[
 \sum_{i,j}|A_{ij}|\le n\|A\|_F\le n\sqrt n\norm{A}=n^{3/2}.
\]
Thus $W\le2n^2+n^{3/2}/h=(2+64/\eps)n^2$. Because $2\eps<1$,
this implies \eqref{eq:weight}. The comparison-floor costs sum to
$O(n^2+\sum|A_{ij}|/h)=O(n^2/\eps)$; all other entrywise work is constant
per entry. These estimates are deterministic.

Let $E=Q-A$. Its entries are independent, centered random variables, each
supported in an interval of length at most $h$. We include the elementary
exponential estimate used here. If $Y$ is centered and supported in an interval
of length $L$, then, for $K(s)=\log\E e^{sY}$, the second derivative is the
variance under an exponentially tilted distribution. Any variable in such an
interval has variance at most $L^2/4$: subtract the midpoint and use
$\operatorname{Var}(Y)\le\E(Y-\mathrm{midpoint})^2$. Hence
$K''(s)\le L^2/4$ and $K(0)=K'(0)=0$, giving
\[
 \E e^{sY}\le e^{s^2L^2/8}.
\]
For fixed unit vectors $u,v\in\R^n$, independence and
$\sum_{ij}u_i^2v_j^2=1$ imply
\[
 \E e^{s u^\top Ev}\le e^{s^2h^2/8}.
\]
Optimizing the exponential Markov bound in $s$ and applying it to both signs
therefore yields
\begin{equation}\label{eq:bilinear}
 \PP\{|u^\top Ev|>t\}\le2e^{-2t^2/h^2}.
\end{equation}

There is a $1/4$-net $\mathcal N$ of the unit sphere with
$|\mathcal N|\le9^n$. Indeed, take a maximal $1/4$-separated set; its disjoint
radius-$1/8$ balls fit in the radius-$9/8$ ball, so volume comparison proves
the cardinality bound, and maximality proves the net property. For every
matrix $E$, approximation of both unit vectors in a maximizing bilinear form
gives
\[
 \norm{E}\le2\max_{u,v\in\mathcal N}|u^\top Ev|.
\]
The two approximation errors are each at most $\norm{E}/4$.
A union bound in \eqref{eq:bilinear}, with $t=\rho/2$, proves
\[
 \PP\{\norm{E}>\rho\}
 \le2\,9^{2n}\exp\!\left(-\frac{\rho^2}{2h^2}\right)
 =2\exp\bigl((2\log9-32)n\bigr).
\]
Since $2\log9<5$, the last quantity is at most $2e^{-27}<10^{-3}$.
The nets and exponential functions are proof devices, not computations
performed by the algorithm.
\end{proof}

\section{Weighted-pattern products with an integer matrix}

We next prove the product construction independently of the random rounding.
An integer pattern is a finite sequence $s=(a_1,\ldots,a_m)$, with
\[
 \omega(s)=\sum_{r=1}^m(1+|a_r|).
\]
The empty pattern has weight zero. Let $\mathcal P_k$ be the collection of all
patterns of weight at most an integer $k\ge1$.

\begin{lemma}[A small catalogue with explicit identifiers]\label{lem:catalogue}
The catalogue satisfies $|\mathcal P_k|\le4^k$. It can be enumerated, together
with a distinct identifier for each pattern and the head and tail identifiers
of every nonempty pattern, in $O(k4^k)$ scalar operations and storage.
Each given pattern in $\mathcal P_k$ can be encoded in $O(k)$ operations.
\end{lemma}

\begin{proof}
Encode an integer coefficient with symbols in $\{0,1,2\}$ as follows:
$0$ is the single symbol $0$; $+r$ is $r$ copies of $1$ followed by $0$;
$-r$ is $r$ copies of $2$ followed by $0$. Concatenate the coefficient codes
and pad the resulting word by symbols $3$ to total length $k$.
The unpadded length is exactly the weight. Zeros terminate coefficients;
the first $3$, if present, starts the padding. Thus the encoding is injective.
There are $4^k$ words of length $k$ over four symbols.

The identifier is the base-four value of the padded word, computed by the
recurrence $\texttt{code}\leftarrow4\,\texttt{code}+\texttt{symbol}$.
Leading zeros cause no ambiguity because every padded word has length $k$.
No bit-packing or digit extraction is used. To enumerate the catalogue,
recursively append all integers of weight at most the remaining budget.
There are at most $4^k$ visited patterns, at most $O(k)$ work at each visit,
and at most $k$ coefficients to store per pattern. Computing the tail code
afresh costs another $O(k)$ per pattern, still within $O(k4^k)$.
The power $4^k$ itself is formed by repeated multiplication.
\end{proof}

\begin{lemma}[Preprocessed integer products]\label{lem:product}
Let $Z\in\Z^{n\times n}$, let $1\le k\le n$, and put
$W=\sum_{ij}(1+|Z_{ij}|)$. In $O(W+k4^k)$ preprocessing operations one can
construct exact product routines for both $Zv$ and $Z^\top v$, valid for
arbitrary real $v$. Each product costs
\begin{equation}\label{eq:productcost}
 O(n4^k+W/k+n)
\end{equation}
operations. The storage, including input, descriptors, and a product's
intermediate table, is $O(n^2+k4^k+n4^k)$ real-scalar words.
\end{lemma}

\begin{proof}
\textit{Packing the rows.}
Give an entry $a$ weight $1+|a|$. If this exceeds $k$, mark it heavy and
store it as a standalone multiplication. In each contiguous run of remaining
light entries, greedily pack entries into maximal consecutive bins whose total
weight is at most $k$. Each bin is a catalogue pattern, recorded by its
identifier and starting column. Its length is at most $k$.

Let $H$ be the number of heavy entries, and let $W_L,W_H$ be the total light
and heavy weights. In a light run, any two consecutive bins have combined
weight greater than $k$: otherwise the first entry of the second bin could
have been added to the first. Pairing consecutive bins shows that the run has
at most $2W_{\rm run}/k+1$ bins. There are at most $H+n$ nonempty light runs.
Since $H\le W_H/k$, the number $S$ of all bins and standalone heavy terms
satisfies
\begin{equation}\label{eq:terms}
 S\le2W_L/k+(H+n)+H\le2W/k+n.
\end{equation}
Each term covers at least one original entry, so also $S\le n^2$.
Scanning costs $O(n^2)$, and encoding a bin costs $O(k)$, including padding.
Thus the scan and encoding cost is
$O(n^2+Sk)=O(W+nk)=O(W)$ because $W\ge n^2$ and $k\le n$.
Heavy coefficients are never expanded in unary.

\textit{All shifted pattern products.}
For a queried vector $v$, build a table $F(j,s)$ for
$j=n,n-1,\ldots,1$ and $s\in\mathcal P_k$. Initialize
$F(n+1,s)=0$ and $F(j,\emptyset)=0$. If $s=(a,s')$ is nonempty, use
\begin{equation}\label{eq:table}
 F(j,s)=a v_j+F(j+1,s').
\end{equation}
Induction on $j$ proves that $F(j,s)$ is the dot product of $s$ with the
segment of $v$ starting at $j$, with zero padding beyond $n$.
In particular it is the exact contribution of a stored bin.
The head and tail identifiers were computed in preprocessing. Each table
entry now costs at most one multiplication and one addition, not $O(k)$
operations. Initializing and filling a direct-address table of
$(n+1)4^k$ slots costs $O(n4^k)$, including unused slots.

For each row, sum one looked-up value for each bin and one scalar product
for each heavy term. This costs $O(S+n)$ operations. Together with
\eqref{eq:terms}, this proves \eqref{eq:productcost}. The table is recomputed
for every queried vector; its entire cost has just been counted.

Apply the same preprocessing separately to $Z^\top$, using the same
catalogue. Its weight is also $W$ and transposing costs $O(n^2)$. Two routines
change the constants only. At most $n^2$ terms are stored per routine; the
catalogue requires $O(k4^k)$ words, and a table requires $O(n4^k)$ words.
The two products can reuse table workspace, proving the storage claim.
\end{proof}

\begin{remark}[Meaning of table access]
All identifiers are integers constructed by the explicitly charged loops and
scalar arithmetic. Their range is $0,\ldots,4^k-1$; ordinary indexed storage
suffices. Alternatively, the identifiers tell a preprocessor which previously
computed scalar to connect to each output sum. Evaluating the resulting
straight-line circuit needs only scalar arithmetic. Neither a hash table nor
a nonuniform advice string about $A$ is required.
\end{remark}

For $Q=hZ$, multiply a routine's output by $h$, at $O(n)$ additional cost.
The same statement applies to $Q^\top$.

\section{A condition-independent approximate-kernel filter}

The next lemma deliberately uses a rectangular matrix. Its nontrivial kernel
is guaranteed by dimension, not by a computed singular vector, a nonsingular
rounded matrix, or consistency of a rounded linear system.

\begin{lemma}[Fixed-time randomized kernel filtering]\label{lem:kernel}
Let $B\in\R^{(d-1)\times d}$ satisfy $\norm{B}^2\le4$, and let $\eta>0$.
Let $\ell$ be the least nonnegative integer with $2^\ell\ge10^9d$.
Let $T$ be the least nonnegative integer satisfying $T\eta^2\ge4\ell$.
Draw $g\sim N(0,I_d)$ and compute
\begin{equation}\label{eq:iteration}
 z_0=g,\qquad z_{t+1}=z_t-\tfrac14B^\top Bz_t,
 \qquad 0\le t<T.
\end{equation}
With probability greater than $0.998$, $z_T\ne0$ and
\begin{equation}\label{eq:approxkernel}
 \norm{Bz_T}\le\eta\norm{z_T}.
\end{equation}
No condition number or spectral gap enters $T$.
\end{lemma}

\begin{proof}
There exists a unit vector $v\in\ker B$; fix one for analysis only.
Put $a=|v^\top g|$ and $G=\norm{g}$. A standard Gaussian has density at most
$1/\sqrt{2\pi}$, so
\[
 \PP\{a<10^{-3}\}\le\sqrt{2/\pi}\,10^{-3}<10^{-3}.
\]
Also $\E G^2=d$, hence Markov's inequality gives
$\PP\{G^2>1000d\}\le10^{-3}$. On their intersection, an event of probability
greater than $0.998$,
\begin{equation}\label{eq:gaussevent}
 a\ge10^{-3},\qquad G^2/a^2\le10^9d.
\end{equation}

Let $H=B^\top B$ and $M=I-H/4$. The eigenvalues $\lambda$ of $H$ lie in
$[0,4]$, so $M$ is a positive semidefinite contraction and $Mv=v$.
The iterate is $z_T=M^{\,T}g$, where $T$ denotes an integer power here,
not transpose. Thus $v^\top z_T=v^\top g$, and
\begin{equation}\label{eq:nonzero}
 \norm{z_T}\ge a>0.
\end{equation}
Use an orthonormal eigenbasis of $H$, with coordinates $g_j$ and
$z_{T,j}=(1-\lambda_j/4)^Tg_j$. The coordinates with
$\lambda_j\le\eta^2/2$ contribute at most
$\frac12\eta^2\norm{z_T}^2$ to
$\norm{Bz_T}^2=\sum_j\lambda_j z_{T,j}^2$.

For the remaining coordinates, the inequality $1-s\le e^{-s}$ for
$0\le s\le1$ gives
\[
 \lambda(1-\lambda/4)^{2T}\le\lambda e^{-T\lambda/2}.
\]
Because $T\ge4/\eta^2$, the function
$\lambda\mapsto\lambda e^{-T\lambda/2}$ is decreasing on
$[\eta^2/2,\infty)$. Consequently the high-eigenvalue contribution is at most
\begin{equation}\label{eq:highpart}
 \frac{\eta^2}{2}\,e^{-T\eta^2/4}G^2.
\end{equation}
By construction, $\ell\ge\log_2(10^9d)\ge\log(10^9d)$ and
$T\eta^2/4\ge\ell$. Therefore \eqref{eq:gaussevent} bounds
\eqref{eq:highpart} by $\eta^2a^2/2$, which is at most
$\eta^2\norm{z_T}^2/2$ by \eqref{eq:nonzero}. Adding the two parts proves
\eqref{eq:approxkernel}. This proof also shows why no additional
$\log(1/\eta)$ factor is needed: the factor $\lambda$ in the residual is
retained when bounding the high-eigenvalue part.
\end{proof}

For our application take $d=n+1$ and
\begin{equation}\label{eq:B}
 B=[\,Q\;\;-c\,].
\end{equation}
On the rounding-success event of Lemma~\ref{lem:round},
\begin{equation}\label{eq:Bnorm}
 \norm{Q}\le1+\rho<17/16,\qquad
 \norm{B}^2\le\norm{Q}^2+\norm{c}^2
 \le(17/16)^2+1=545/256<4.
\end{equation}
For $z_t=(u_t,\alpha_t)$, a step of \eqref{eq:iteration} is exactly
\begin{equation}\label{eq:implementedstep}
 \begin{aligned}
 r_t&=Qu_t-c\alpha_t,\\
 u_{t+1}&=u_t-\tfrac14Q^\top r_t,\\
 \alpha_{t+1}&=\alpha_t+\tfrac14c^\top r_t.
 \end{aligned}
\end{equation}
It uses one product with each of $Q,Q^\top$ and $O(n)$ further operations.
In particular neither $B^\top B$ nor a dense Gram matrix is formed.
The vector $v$ in the proof is not computed.

\section{Conversion to backward error with \texorpdfstring{$b$}{b} unchanged}

\begin{lemma}[Clamped projective conversion]\label{lem:convert}
Assume \eqref{eq:parameters}, $\norm{Q-A}\le\rho$, and
$z=(u,\alpha)\ne0$ satisfies
\begin{equation}\label{eq:conversionhyp}
 \norm{Qu-c\alpha}\le\eta\norm{z}.
\end{equation}
Then $u\ne0$. Define $\sgnp(\alpha)=1$ for $\alpha\ge0$ and $-1$ otherwise,
and put
\begin{equation}\label{eq:clamp}
 \alpha'=\sgnp(\alpha)\max\{\abs{\alpha},\eta\norm{u}\},
 \qquad x=\frac{\beta}{\alpha'}u.
\end{equation}
This vector is nonzero and satisfies
\begin{equation}\label{eq:converted}
 \frac{\norm{Ax-b}}{\norm{x}}\le\rho+4\eta=5\eps/8,
 \qquad \norm{x}\le\beta/\eta.
\end{equation}
In particular the definition remains valid when $\alpha=0$.
\end{lemma}

\begin{proof}
If $u=0$, then $\alpha\ne0$ and
$\norm{Qu-c\alpha}=|\alpha|=\norm{z}$, contradicting $\eta<1$.
Moreover,
\[
 |\alpha|=\norm{c\alpha}
 \le\norm{Qu}+\norm{Qu-c\alpha}
 \le(1+\rho)\norm{u}+\eta\norm{z}.
\]
Using $\norm{z}\le\norm{u}+|\alpha|$ gives
\begin{equation}\label{eq:projectiveratio}
 \frac{\norm{z}}{\norm{u}}
 \le\frac{2+\rho}{1-\eta}<3,
\end{equation}
since $\rho=\eta<1/16$. The clamping rule gives
$|\alpha'|\ge\eta\norm{u}>0$ and
$|\alpha'-\alpha|\le\eta\norm{u}$, including when $\alpha=0$.
Thus $x\ne0$ and $\norm{x}\le\beta/\eta$. Since $b=\beta c$ exactly,
\begin{align*}
 \frac{\norm{Ax-b}}{\norm{x}}
 &=\frac{\norm{Au-c\alpha'}}{\norm{u}}\\
 &\le\norm{A-Q}
   +\frac{\norm{Qu-c\alpha}}{\norm{u}}
   +\frac{|\alpha'-\alpha|}{\norm{u}}\\
 &\le\rho+\eta\frac{\norm{z}}{\norm{u}}+\eta
 \le\rho+4\eta.
\end{align*}
This establishes the claimed guarantee for the original $A$ and $b$.
\end{proof}

For completeness, the backward perturbation is fully explicit. Given any
nonzero output $x$, set
\begin{equation}\label{eq:rankone}
 \Delta A=\frac{(b-Ax)x^\top}{\norm{x}^2}.
\end{equation}
Then $(A+\Delta A)x=b$ and
$\norm{\Delta A}=\norm{Ax-b}/\norm{x}$. Conversely,
$\Delta A x=b-Ax$ implies
$\norm{\Delta A}\ge\norm{Ax-b}/\norm{x}$, so this is exactly the minimum
spectral-norm perturbation, not merely an upper bound on a different metric.
On success, \eqref{eq:converted} and $\norm{A}=1$ give
$\norm{\Delta A}/\norm{A}\le5\eps/8$. No perturbation of $b$ is needed.

\clearpage
\section{Complete algorithm and deterministic accounting}

The following is the exact-real algorithm. The shorthand for rounding and
product preprocessing denotes precisely the loops proved above.
The iteration count is implemented by a comparison, rather than a ceiling
primitive.

\begin{lstlisting}[caption={Exact-real solver; all loops have deterministic cost bounds.}]
Input: entries of A, b; promised ||A||_2 = 1; b != 0; 0 < eps < 1/2
beta = sqrt(sum_i b[i]^2); c = b / beta
rho = eta = eps / 8; h = eps / (64 * sqrt(n))

For each entry A[i,j]:
    t = A[i,j] / h
    compute m = floor(t) by the comparison/increment loops
    draw three fresh independent standard Gaussians
    obtain U from (1 + G1 / sqrt(G1^2 + G2^2 + G3^2)) / 2
        using U = 1/2 when the square-root denominator is zero
    Z[i,j] = m + 1 if U < t - m, and m otherwise

k = 1; power = 16
while 16 * power <= n:
    power = 16 * power; k = k + 1
preprocess weighted-pattern products for Z and transpose(Z)
Q_product(v)  = h * Z_product(v)
QT_product(v) = h * ZT_product(v)

ell = 0; power = 1
while power < 10^9 * (n + 1):
    power = 2 * power; ell = ell + 1
Draw n + 1 fresh independent standard Gaussians as (u, alpha)
t = 0
while t * eta^2 < 4 * ell:
    r = Q_product(u) - c * alpha
    u = u - QT_product(r) / 4
    alpha = alpha + dot(c, r) / 4
    t = t + 1

s = sqrt(dot(u,u))
if s == 0: return beta * e_1
alpha_prime = max(abs(alpha), eta * s)
if alpha < 0: alpha_prime = -alpha_prime
return (beta / alpha_prime) * u
\end{lstlisting}

\subsection{A uniform bound valid even for small \texorpdfstring{$n$}{n}}

\begin{lemma}[Cancellation parameters]\label{lem:budget}
For the $k,\ell,T$ used in the algorithm,
\begin{equation}\label{eq:kproperties}
 1\le k\le n,\qquad n<16^{k+1},\qquad k4^k\le4n,
\end{equation}
\begin{equation}\label{eq:timeproperties}
 \ell\le4k+35\le39k,\qquad T\le10000\,k/\eps^2.
\end{equation}
They can be determined with the permitted arithmetic and comparisons.
\end{lemma}

\begin{proof}
For $n\ge16$, $16^k\le n<16^{k+1}$. Since $k\le4^k$ for positive integers,
$k4^k\le16^k\le n$. For $1\le n<16$, $k=1$ and $k4^k=4\le4n$.
The remaining assertions in \eqref{eq:kproperties} follow directly.
Since $10^9<2^{30}$ and $n+1\le2n<2\,16^{k+1}$,
\[
 10^9(n+1)<2^{4k+35},
\]
so the least possible $\ell$ is at most $4k+35\le39k$.
The least $T$ with $T\eta^2\ge4\ell$ obeys
\[
 T\le\frac{4\ell}{\eta^2}+1
 =\frac{256\ell}{\eps^2}+1
 \le\frac{9984k}{\eps^2}+1
 \le\frac{10000k}{\eps^2}.
\]
Repeated multiplication and comparison implement the $k$ and $\ell$ loops;
the iteration loop realizes exactly this $T$. Logarithms, floors, and ceilings
in their mathematical descriptions are not unit-cost operations in the
algorithm.
\end{proof}

\begin{proof}[Proof of Theorem~\ref{thm:main}]
\textit{Correctness and probability.}
By Lemma~\ref{lem:round}, the rounding-success event
$\mathcal E=\{\norm{Q-A}\le\rho\}$ has probability greater than $0.999$.
Conditional on any realized $Q$ in this event, \eqref{eq:Bnorm} holds, and
the fresh Gaussian start permits application of Lemma~\ref{lem:kernel}.
With conditional probability greater than $0.998$ the computed vector is
nonzero and satisfies \eqref{eq:conversionhyp}. Lemma~\ref{lem:convert}
then gives the claimed $5\eps/8$ error for the original input.
A union bound bounds the total failure probability by a quantity strictly
less than $0.001+0.002=0.003$. The conditional argument does not require
independence of the success events; only the fresh starting draw is needed.

\textit{All realizations are well-defined.}
The Gaussian-to-uniform construction has an explicit branch at its zero
denominator. Also $\beta>0$, $h>0$, and no division takes place inside the
filter except by $4$. If $\norm{u}=0$ at the end, the fallback $\beta e_1$
is nonzero. Otherwise the clamped denominator is nonzero. Consequently every
realization, including those outside the success event, returns a nonzero
vector after the same bounded number of iterations. Clamping gives
$\norm{x}\le\beta/\eta$; the fallback also satisfies this bound because
$\eta<1$. There is no condition-number-sensitive stopping rule.

\textit{Operation count.}
Rounding costs $O(n^2/\eps)$. By Lemma~\ref{lem:product}, both product
routines cost $O(W+k4^k)$ to preprocess. One filter step costs
$O(n4^k+W/k+n)$, including both product tables and all updates.
Normalization, loop setup, and final conversion cost at most $O(n+k+\ell)$
in addition. The iteration comparison and counter cost $O(T)$ in total and
are absorbed by the step bound. Therefore the entire cost is at most
\begin{align}
 &O(n^2/\eps)+O(W+k4^k)
   +O\!\left(\frac{k}{\eps^2}
                  \left(n4^k+\frac Wk+n\right)\right) \notag\\
 &\hspace{18mm}=O(n^2/\eps)
   +O\!\left(\frac{kn4^k+W+nk}{\eps^2}\right)
   =O(n^2/\eps^3).\label{eq:finalcost}
\end{align}
The last equality uses $k4^k\le4n$, $k\le n$, and
$W\le65n^2/\eps$. Every estimate is uniform in all inputs and all random
outcomes. This proves existence of an absolute operation-count constant $C$
with $q=3$.

\textit{Storage.}
The input and two term lists require $O(n^2)$ words.
By \eqref{eq:kproperties}, $k4^k=O(n)$ and $n4^k=O(n^2)$, so the catalogue
and reusable product tables fit in $O(n^2)$ words. The vectors, counters,
and temporary scalars need less space. This proves the storage claim.
\end{proof}

An optional early exit on
$\norm{Qu-c\alpha}\le\eta\sqrt{\norm{u}^2+\alpha^2}$ is valid in exact
arithmetic and does not increase the worst-case cost. It tests the augmented
residual using a preprocessed product. Recomputing the original dense residual
$Ax-b$ on every iteration is \emph{not} part of the algorithm; that would
reintroduce an avoidable logarithmic cost. A single final original-system
check or construction of \eqref{eq:rankone} costs $O(n^2)$ and is harmless.

\section{Reproducibility and scope}

The accompanying archive contains the mathematical source, a PDF, exact
component tests, a floating-point solver prototype, reproducible experiments,
and a separate audit of the model and edge cases. These serve different roles.

\subsection{Exact component checks}

The standard-library test suite uses integers and \texttt{fractions.Fraction}
for exact checks. All nine test groups pass. They cover injective pattern
identifiers, all shifted pattern products, exhaustive small greedy packings,
exact forward and transpose products, comparison-based floors, the uniform
parameter inequalities, rational diagonal filtering examples, the
$\alpha=0$ and tiny-$\alpha$ clamping cases, and the error-budget constants.
The tests include $147$ rational matrix fixtures with both product directions,
packing budgets $1$ through $6$, and integer dimension boundary checks near
$16^j$ for $j\le100$. These are implementation checks, not a proof of the
universal theorem; the proof is Sections 2--6.

\subsection{Floating-point illustrations}

The NumPy prototype is explicitly not a finite-precision theorem. It is tested
on $15$ fixed-seed systems, including ill-conditioned, nonnormal, and one
singular inconsistent example. All final original-system checks passed, and
all reported errors are below $5\eps/8$. Table~\ref{tab:tests} gives selected
runs; the complete data, seeds, and version metadata are in the archive.

\begin{table}[ht]
\centering
\small
\begin{tabular}{@{}lrrrr@{}}
\toprule
Test family & $n$ & $\eps$ & Steps & $\norm{Ax-b}/\norm{x}$\\
\midrule
Diagonal, smallest entry $10^{-250}$ & 4 & 0.25 & 171 & $3.1092\!\times\!10^{-2}$\\
Small singular-value outlier & 32 & 0.10 & 240 & $1.4401\!\times\!10^{-2}$\\
Geometric spectrum & 64 & 0.25 & 75 & $4.1896\!\times\!10^{-2}$\\
Nonnormal bidiagonal & 24 & 0.25 & 424 & $3.1829\!\times\!10^{-2}$\\
Singular, inconsistent (extension) & 16 & 0.25 & 51 & $4.4016\!\times\!10^{-2}$\\
\bottomrule
\end{tabular}
\caption{Illustrations with augmented-residual early exit. Main experiment
runs use dense products for convenience; separate same-draw comparisons test
the compressed implementation. The table is not a speed benchmark.}
\label{tab:tests}
\end{table}

The dense and compressed backends were compared end to end on five of the
same draws, with maximum observed relative output difference below
$2.0\times10^{-16}$. Separate product comparisons exercise budgets
$k=2,\ldots,7$, rather than only the small default budgets at modest $n$.
In $300$ diagnostic rounding trials the largest observed ratio
$\norm{Q-A}/(\eps/8)$ was approximately $0.1186$. Diagnostic spectral norms
and matrix-family construction occur outside the proved algorithm.
None of these finite samples establishes its universal success probability.

The solver prototype uses floating-point Gaussians and arithmetic, and uses a
Python ceiling for its diagnostic step cap; the paper gives comparison-based
implementations for every restricted-model operation. An optional NumPy floor
shortcut, optional dense products, and an optional user-imposed iteration cap
are clearly marked in the code. A truncated run does not inherit the
fixed-time theorem. Exact circuit tests require only the Python standard
library; the numerical experiments require NumPy.

\subsection{What the result does and does not establish}

The answer to the displayed IE-12 cost question is affirmative with an
explicit exponent $q=3$. The success criterion is precisely $A$-only
normwise backward error for the original data. No condition-number bound,
input-distribution assumption, rounded-matrix invertibility, or restriction
on the direction or magnitude of $b$ is used.

The method still takes $O(\eps^{-2}\log(n+1))$ filter steps. It removes the
dimension-dependent logarithm from \emph{total arithmetic cost} through
entry-dependent preprocessing, rather than proving a dimension-free black-box
iteration rate. Constants are deliberately generous and no optimal exponent
in $\eps$ is claimed. At modest dimensions the chosen catalogue budget is
small, so this construction should not be interpreted as a demonstrated
practical improvement over established solvers.

The theorem counts exact-real scalar operations and real-scalar words.
It does not bound finite-precision roundoff or the bit lengths of exact
intermediates. In particular, a small backward error need not imply a small
forward error for an ill-conditioned input. The explicit perturbation
\eqref{eq:rankone}, not closeness to $A^{-1}b$, is the certificate.
The result is presented with a self-contained proof; independent mathematical
review would be appropriate before treating it as an externally established
resolution of the repository entry.

\clearpage
\begin{thebibliography}{9}
\bibitem{ie12}
\emph{OpenProblemsInNLA}, IE-12, ``Near-quadratic solution cost at a prescribed
backward error,'' entry marked last checked 10 September 2026. Website read
12 September 2026.\\
\url{https://github.com/ajt60gaibb/OpenProblemsInNLA/blob/main/linear-systems-and-elimination/IE-12/README.md}

\bibitem{dnr}
M. Derezi\'nski, Y. Nakatsukasa, and E. Rebrova,
\emph{Towards Universal Convergence of Backward Error in Linear System
Solvers}, arXiv:2604.16075v2, 22 May 2026. In particular, Corollaries 17 and 25
and Section 7.\\
\url{https://arxiv.org/html/2604.16075v2}

\bibitem{mailman}
E. Liberty and S. W. Zucker,
``The Mailman algorithm: A note on matrix--vector multiplication,''
\emph{Information Processing Letters} \textbf{109}(3), 179--182 (2009).
DOI: 10.1016/j.ipl.2008.09.028. Author manuscript: \\
\url{https://www.cs.yale.edu/homes/el327/papers/mailmanAlgorithm.pdf}
\end{thebibliography}

\end{document}
